\begin{center}\begin{tikzpicture}[font=\small,>=stealth,line join=round]
\draw[gray!40,rounded corners=3pt] (0,0) rectangle (14.2,8.7);\node at (7.1,8.1) {例8：根号差为零，就是到两定点一样远};
\begin{scope}[shift={(2.5,2.2)},x=1.5cm,y=1.5cm]\draw[->,gray] (-0.6,0)--(2.6,0) node[right]{$x$};\draw[->,gray] (0,-.35)--(0,2.8) node[left]{$y$};\draw[orange!85!black,line width=1.6pt] (0,2)--(0.25,0);\draw[violet,thick] (2,1)--(0.25,0);\draw[teal!80!black,thick] (0,2)--(2,1);\fill[teal!80!black] (0,2) circle(2.2pt) node[above left=4pt]{$A$};\fill[teal!80!black] (2,1) circle(2.2pt) node[above right=4pt]{$B$};\fill[blue] (0.25,0) circle(2.2pt) node[below=4pt]{$P$};\end{scope}\begin{scope}[shift={(2.5,2.2)},x=1.5cm,y=1.5cm]\draw[gray,dashed] (.1,-.3)--(1.55,2.6);\fill[gray] (1,1.5) circle(1.8pt);\node[above left=7pt] at (1,1.5){$M$};\draw[gray] (1.134,1.433)--(1.201,1.567)--(1.067,1.634);\end{scope}\node[anchor=west,align=left] at (8.5,4.4) {$PA=PB$\\[12pt]$P$ 在 $AB$ 的垂直平分线上\\[12pt]与 $x$ 轴交于 $x=1/4$\\[12pt]有理化后分子为 $4x-1$\\[10pt]所以零点仅为 $1/4$。};
\node[text width=12.8cm,align=center] at (7.1,.65) {橙、紫两段等长；灰点 $M$ 为 $AB$ 中点，灰虚线为垂直平分线。};
\end{tikzpicture}\end{center}
